\documentclass[14pt, a4paper]{extarticle}
\usepackage[T2A]{fontenc}
\usepackage[utf8]{inputenc}
\usepackage[english, russian]{babel}
\usepackage[left=30mm, right=10mm, top=20mm, bottom=25mm]{geometry}

\usepackage{misccorr} % в заголовках появляется точка, но при ссылке на них ее нет
\usepackage{indentfirst} % после заголовков ставится абзацный отступ

% Заголовки section по центру %
\makeatletter
\renewcommand\section{\@startsection {section}{1}{\z@}%
                                   {-3.5ex \@plus -1ex \@minus -.2ex}%
                                   {2.3ex \@plus.2ex}%
                                   {\centering\normalfont\Large\bfseries}}
\makeatother
%%%%%

% Точки в списке литературы %
\makeatletter
\renewcommand\@biblabel[1]{#1.}
\makeatother
%%%%%

\begin{document}
\thispagestyle{empty}
{\scriptsize
\begin{center}
ФЕДЕРАЛЬНОЕ ГОСУДАРСТВЕННОЕ ОБРАЗОВАТЕЛЬНОЕ БЮДЖЕТНОЕ УЧРЕЖДЕНИЕ\\ 
ВЫСШЕГО ОБРАЗОВАНИЯ    
\end{center}
}
{\bf
\begin{center}
РГАУ-МСХА имени К.А. Тимирязева\\ 
\end{center}
}
\hrule
\bigskip

{\small 
\begin{center}
\textbf{
Кафедра <<Статистики и кибернетики>>}
\end{center}
}
\bigskip

Профиль: Информационные системы и технологии

\medskip
Направленность: Компьютерные науки и технологии искусственного интеллекта

\vskip 2cm

\begin{center}
\textbf{КУРСОВОЙ ПРОЕКТ}
\end{center}

\begin{center}
по дисциплине:\\ 
<<Алгоритмизация и программирование>>    
\end{center}


\bigskip
\begin{center}
ТЕМА:\\ 
Название темы
\end{center}


\vskip 4cm

\hskip .4\textwidth Преподаватель:

\hskip .4\textwidth Студент:

\hskip .4\textwidth Курс: первый

\vskip 7.0cm

\begin{center}
Липецк 2026
\end{center}

\newpage
\tableofcontents
\newpage
\addcontentsline{toc}{section}{ВВЕДЕНИЕ}
\section*{ВВЕДЕНИЕ}

\newpage
\section{Глава 1. Теоретическая часть}
Text \cite{book1}.

Текст текст \( \int_a^b f(x)dx. \)
\[
\int_a^b f(x)dx.
\]
\subsection{Пункт}
\begin{equation}\label{integral}
    \int_{a}^b \sin x dx
\end{equation}
\begin{equation}\label{summa}
    \sum\limits_{i=1}^n x_i
\end{equation}
\begin{equation}\label{alphabeta}
    \alpha +\beta
\end{equation}
\subsection{Пункт}

\newpage
\section{Глава 2. Практическая часть}
text (\ref{summa})
\subsection{Пункт}
text (\ref{alphabeta})
\subsection{Пункт}

\newpage
\section{НАЗВАНИЕ}
\subsection{Пункт}
\subsection{Пункт}

\newpage
\section{НАЗВАНИЕ}
\subsection{Пункт}
\subsection{Пункт}

\newpage
\addcontentsline{toc}{section}{ЗАКЛЮЧЕНИЕ}
\section*{ЗАКЛЮЧЕНИЕ}

\newpage
\renewcommand\refname{\centering СПИСОК ИСПОЛЬЗОВАННЫХ ИСТОЧНИКОВ}
\clearpage% or cleardoublepage
\addcontentsline{toc}{section}{СПИСОК ИСПОЛЬЗОВАННЫХ ИСТОЧНИКОВ}
\begin{thebibliography}{99}
\bibitem{book3} Author Book2

\bibitem{book1} Author Book1

\bibitem{book2} Author Book2
\end{thebibliography}

\newpage
\addcontentsline{toc}{section}{ПРИЛОЖЕНИЯ}
\section*{ПРИЛОЖЕНИЯ}

\end{document}
